\input ../common.tex

\graphicspath{{fig/}{../fig/}{../01/fig/}}
\usepackage[export]{adjustbox}

\institute[Nick Papior; DTU Compute]{\begin{tikzpicture}
      \node[shape=rectangle split,rectangle split parts=2,anchor=base] at (0,0)
      {SIESTA School};
    \end{tikzpicture}}

\date{15 November 2024}
\title{Intro: Non-equilibrium Green function theory}
\author{Nick Papior}


\begin{document}

\begin{frame}
  \titlepage
\end{frame}

\begin{frame}
  \frametitle{Outline}
  \tableofcontents
\end{frame}

\showbondcurrentfalse
\showdensitymatrixweighingfalse
\showintromotivationfalse
\showintrotbtnterminalfalse
\shownegftransmissionfalse

\inputpath[../01]{introduction.tex}
\inputpath[../01]{se.tex}
\inputpath[../01]{green.tex}

\begin{frame}

  \vfill

  \centering
  \begin{minipage}{.7\linewidth}

  \begin{block}{Important}

    Understanding self-energies is like, really, really important.

  \end{block}

  \end{minipage}

  \vfill
\end{frame}

\inputpath[../01]{se.tex}
\inputpath[../01]{negf.tex}
%\inputpath[../03]{contour.tex}

\inputpath[../03]{steps.tex}
%\inputpath[../03]{summary.tex}


\section{Tutorial}
\begin{frame}
  \frametitle{Tutorial}

  \begin{itemize}
    \item Tutorial focuses on \emph{interpreting physical} things, such as potential drop, projected
      density of states

    \item More tutorials are available here:

      \url{https://github.com/zerothi/ts-tbt-sisl-tutorial}

    \item Doing this tutorial can be hard on Marenostrum (requires copying data back to
      your PC).

      Please use Google colab or
      your own computer.

    \item On the Discord channel in Day 5, you will find the link to the tutorial.

  \end{itemize}

\end{frame}

\end{document}
